\section{物理地址、兼容地址图与复位坐标}

处理器访问物理地址时，平台接收的不只是一个整数。事务还带有读写方向、访问
宽度、缓存属性和权限来源；响应方返回数据或错误。真实总线可能把一次请求拆成
地址阶段与数据阶段，并用等待状态处理慢设备。QEMU 在软件中实现同样的可见
结果，不要求模拟每根电气引脚。

\begin{pdfdiagram}
  \node[os block] (request) {CPU 请求\\地址、读写、宽度、属性};
  \node[os state,right=of request] (fabric) {平台互连\\范围比较与路由};
  \node[os hardware,above right=5mm and 12mm of fabric] (memory) {DRAM/ROM\\返回数据};
  \node[os hardware,below right=5mm and 12mm of fabric] (device) {设备寄存器\\改变状态};
  \node[os state,right=30mm of fabric] (response) {完成响应\\数据或错误};
  \draw[os arrow] (request) -- (fabric);
  \draw[os arrow] (fabric) -- (memory);
  \draw[os arrow] (fabric) -- (device);
  \draw[os arrow] (memory) -- (response);
  \draw[os arrow] (device) -- (response);
\end{pdfdiagram}
\diagramnote{地址译码只决定响应者；访问语义还由宽度、顺序、缓存属性和设备状态共同决定。}

访问 RAM 通常可以重复读取而不改变内容，访问设备状态寄存器却可能具有读清除
副作用；写 RAM 表示保存数据，写命令寄存器可能启动异步操作。把 MMIO 强制
转换为普通指针后，C++ 类型系统不会自动知道这些区别，驱动必须在接口层保留
协议语义。

\topic{PC 物理地址图由兼容性层层叠加}
最初 8086 只有 1 MiB 空间。低端区域用于中断向量与 RAM，高端保留显存、
扩展 ROM 和系统固件。处理器地址扩展到 32 位和 64 位后，旧设备窗口仍需
兼容，PCI MMIO、APIC 和固件映射继续占据若干范围。现代 PC 的物理地址空间
因此不是一整块连续 DRAM。

\begin{osgridlongtable}{L{0.24\textwidth}L{0.26\textwidth}L{0.36\textwidth}}
\textbf{示意范围} & \textbf{传统用途} & \textbf{本项目使用原则}\\ \hline
\endfirsthead
\textbf{示意范围} & \textbf{传统用途} & \textbf{本项目使用原则}\\ \hline
\endhead
\code{0x00000000} 起 & 低端 RAM、向量和固件数据 & 明确保留后再放栈与启动结构\\ \hline
\code{0x000A0000} 附近 & VGA 与适配器窗口 & 不当作普通可用 RAM\\ \hline
\code{0x00100000} 以上 & 扩展 RAM & 后续装载内核和页表\\ \hline
4 GiB 以下高端 & PCI 窗口与平台设备 & 依据机器模型排除\\ \hline
\code{0xFFFE0000} 起 & 本项目 128 KiB ROM & 只读执行与常量区域\\ \hline
\end{osgridlongtable}

表中除项目 ROM 外是解释兼容结构的示意，不是可以硬编码的完整内存图。后续
内核必须读取固件交付的内存区间并验证重叠，而不是假定“内存大小”等于从零
开始的连续 RAM。

\topic{复位取指的三个坐标系逐项对应}
链接器中的地址是符号运行地址，ROM 文件中的位置是文件偏移，处理器发出的是
物理地址。128 KiB 等于 \code{0x20000}，ROM 顶端对齐到
\code{0x100000000}，所以起点为：
\[
  0x100000000-0x20000=0xFFFE0000.
\]
复位物理地址 \code{0xFFFFFFF0} 对应文件偏移：
\[
  0xFFFFFFF0-0xFFFE0000=0x1FFF0.
\]
入口 \code{0xFFFFF000} 对应 \code{0x1F000}。这两个偏移相差
\code{0xFF0}，复位位置又距 ROM 末端只有 16 字节。

\begin{osgridlongtable}{L{0.22\textwidth}L{0.22\textwidth}L{0.22\textwidth}L{0.22\textwidth}}
\textbf{对象} & \textbf{物理地址} & \textbf{文件偏移} & \textbf{作用}\\ \hline
\endfirsthead
\textbf{对象} & \textbf{物理地址} & \textbf{文件偏移} & \textbf{作用}\\ \hline
\endhead
ROM 起点 & \code{FFFE0000} & \code{00000} & 镜像映射基准\\ \hline
入口 & \code{FFFFF000} & \code{1F000} & 完整初始化代码\\ \hline
复位向量 & \code{FFFFFFF0} & \code{1FFF0} & 三字节近跳转\\ \hline
ROM 末端 & \code{FFFFFFFF} & \code{1FFFF} & 最后一个可寻址字节\\ \hline
\end{osgridlongtable}

\topic{相对跳转依赖下一指令地址}
x86 相对位移以完整指令之后的地址为基准。复位跳转长度
为 3，因此基准 IP 为 \code{0xFFF3}。16 位位移 \code{0xF00D}
符号扩展为 \(-4083\)，结果为 \code{0xF000}。机器码按小端序保存为
\code{E9 0D F0}。

如果误用 \code{target-current}，会多出三个字节；如果把 \code{F0 0D}
按阅读顺序写入镜像，位移变成另一数值；如果使用远跳转，CS 隐藏基址会重新
装载。三类错误都能生成“看起来像跳转”的字节，却不会到达预期物理地址，所以
审计器必须实际解码最终镜像。

\section{Port I/O、设备寄存器与故障诊断}

\code{IN} 与 \code{OUT} 使用 16 位端口号，不经过普通内存页表翻译。
立即数形式只能编码 8 位端口，COM1 的 \code{0x3F8} 超出范围，因此代码先
把端口号装入 DX，再用 \code{out dx,al}。数据宽度由 AL、AX 或 EAX
决定，不由端口号决定。

\begin{osgridlongtable}{L{0.20\textwidth}L{0.24\textwidth}L{0.24\textwidth}L{0.20\textwidth}}
\textbf{指令} & \textbf{端口来源} & \textbf{数据宽度} & \textbf{方向}\\ \hline
\endfirsthead
\textbf{指令} & \textbf{端口来源} & \textbf{数据宽度} & \textbf{方向}\\ \hline
\endhead
\code{in al,dx} & DX 低 16 位 & 8 位 & 设备到 CPU\\ \hline
\code{in ax,dx} & DX 低 16 位 & 16 位 & 设备到 CPU\\ \hline
\code{out dx,al} & DX 低 16 位 & 8 位 & CPU 到设备\\ \hline
\code{out dx,ax} & DX 低 16 位 & 16 位 & CPU 到设备\\ \hline
\end{osgridlongtable}

IDE 数据端口需要 16 位读取，UART 寄存器需要 8 位访问。统一包装端口号不能
抹去寄存器宽度，驱动接口应让错误宽度难以表达。

\topic{设备寄存器不是可随意重排的变量}
编译器可能重排普通内存访问，处理器和总线也可能缓冲写入。端口 I/O 指令具有
x86 规定的顺序属性，但 MMIO 设备还需要合适的缓存类型和屏障。驱动必须依据
设备规范区分配置写入、命令提交、状态读取和完成确认。

UART 的 DLAB 展示了最小状态依赖：同一个端口偏移在 DLAB 为 0 和 1 时指向
不同寄存器。若中断或另一个执行流在配置中途改变 LCR，下一次访问的语义就会
变化。早期固件通过关闭中断和单线程执行避免并发，成熟驱动则要用锁保护整个
寄存器事务。

\topic{ROM、RAM 与设备失败具有不同诊断形态}
ROM 地址错误常表现为执行填充值，RAM 地址错误可能破坏栈后延迟爆发，设备
地址错误则可能读到固定全一或访问另一个设备。三者都可能最终表现为“串口没有
输出”，但定位路径不同。

调试先检查 QEMU 参数是否装入正确镜像，再检查复位偏移和反汇编；随后确认栈
写入落在 RAM，最后跟踪端口事务与状态位。按事务层次收缩范围，比反复修改
输出函数更快找到根因。
